Repozitář a video tutoriály (přístup a provoz)
Online repozitář slouží jako „living document“ s šablonami, video‑návody a verzovanými materiály pro rychlé nasazení.
Přístup a základní orientace (krátký návod)
\begin{itemize}
\item Kde: centrální vstup ai.vut.cz; doporučená struktura: \texttt{/prompts}, \texttt{/templates}, \texttt{/videos}, \texttt{/case‑studies}, \texttt{/blueprints}.
\item Rychlý start: stáhněte ZIP nebo klonujte, prostudujte \texttt{README.md} a otevřete \texttt{/videos} s titulky a časovými razítky.
\end{itemize}
Video tutoriály — doporučené postupy pro učitele
\begin{itemize}
\item Každé video má abstrakt, délku a klíčové časové body.
\item Témata zahrnují praktické návody (např. Minecraft Education) \cite{kb_ai_101_tipu_pdf_04e4a115_page_15_2} a integrace Microsoft (OneNote, Copilot, živé titulky) \cite{kb_ai_101_tipu_pdf_04e4a115_page_8_1, kb_ai_101_tipu_pdf_04e4a115_page_16_1}.
\end{itemize}
Mechanismus kvartálních aktualizací a sběru zpětné vazby
\begin{itemize}
\item Příspěvky průběžně; formální vydání kvartálně. Použijte issue šablonu nebo pull request s popisem pedagogického dopadu a testováním.
\item Ediční tým kontroluje etiku, GDPR a kvalitu; každé vydání obsahuje release notes.
\end{itemize}
Podpora a eskalace
\begin{itemize}
\item Pro dotazy a chyby použijte issues nebo formulář „Kontakt / Podpora“; technické/institucionální otázky směřujte na lokální IT/metodické kontakty (kap. 9.7).
\end{itemize}
Poznámka o ověřeném obsahu
\begin{itemize}
\item Videa a šablony odkazují na ověřené příklady a v popisu vždy uveďte doporučené licence a omezení z hlediska ochrany osobních údajů \cite{kb_ai_101_tipu_pdf_04e4a115_page_15_2, kb_ai_101_tipu_pdf_04e4a115_page_8_1, kb_ai_101_tipu_pdf_04e4a115_page_16_1}.
\end{itemize}